\documentclass[11pt]{article}

\usepackage{fullpage}
\usepackage{amsmath, amssymb, bm, cite, epsfig, psfrag}
\usepackage{graphicx}
\usepackage{float}
\usepackage{amsthm}
\usepackage{amsfonts}
\usepackage{cite}
\usepackage{hyperref}
\usepackage{pgf,tikz}
\usepackage{enumitem}
\usepackage{mathtools}
\usepackage{siunitx}
\usepackage{../../styles/course}


\begin{document}

\title{Introduction to Hardware Design\\
Basic Digital Logic:  Figures and Code Snippets}
\author{Profs. Sundeep Rangan, Siddharth Garg}
\date{}

\maketitle

\subsection{ ReLU with a linear input}

We consider implementing the ReLU function:
  \[
    y = \max\{ ax + b, 0\}
  \]
  where $x$, $a$, and $b$ are integer inputs.  This function arises commonly
  in machine learning.
  We first consider a purely combinational implementation of this function.
  A possible SystemVerilog implementation is as follows:

\begin{systemverilogcode}
    module relu_lin (
      input logic signed [31:0] x,
      input logic signed [31:0] a,
      input logic signed [31:0] b,
      output logic signed [31:0] y
    );
      always_comb begin
        logic signed [31:0] mult_out, add_out;
        mult_out = x * a;
        add_out = mult_out + b;
        y = (add_out > 0) ? add_out : 0;
      end
    endmodule
  \end{systemverilogcode}

  When this module is synthesized, a potential block diagram is as shown in
  Figure~\ref{fig:relu_comb}.

  
\begin{figure}
\centering

  \begin{tikzpicture}[node distance=1cm]

    % x, a and multipliers
    \node (x) {\texttt{x}};
    \node [circop, right=1cm of x,fill=green!20] (mul) {$\times$};
    \node [above of=mul, yshift=0.5cm] (a) {\texttt{a}};
    \draw[->] (x) -- (mul);
    \draw[->] (a) -- (mul);

    % b and adder
    \node[circop,right=2cm of mul,fill=green!20] (add) {$+$};
    \node [above of=add, yshift=0.5cm] (b) {\texttt{b}};
    \draw[->] (mul) -- node[midway,above] {\texttt{mult\_out}} (add);
    \draw[->] (b) -- (add);

    % max block
    \node[draw, right=2cm of add,minimum height=1cm,fill=green!10] (max) {$\max\{u,0\}$};
    \draw[->] (add) -- node[midway,above] {\texttt{add\_out}} (max);
    \draw[->] (max) -- ++(1.5,0) node[right] {\texttt{y}};
\end{tikzpicture}
  \caption{Fully combinational block diagram for ReLU with a linear input.}
  \label{fig:relu_comb}
\end{figure}

\subsection{Two input circuit}
Next, we consider a simple two input circuit with the following SystemVerilog code:
\begin{systemverilogcode}
xsq = x * x;
r = (x > 0) ? x : 0;
y = xsq + r;
\end{systemverilogcode}

  
\begin{figure}
\centering

  \begin{tikzpicture}[node distance=1cm]

    % x, a and multipliers
    \node (x) {\texttt{x}};
    \node [circop, right=1cm of x,fill=green!20] (mul) {$\times$};

    \coordinate[right=0.3cm of x] (xsplit);
    \draw[->] (x) -- (xsplit);
    \draw[->] (xsplit) -- (mul);   

    % max block
    \node[draw, right=1cm of x,minimum height=1cm,fill=green!10,yshift=-1.5cm] 
      (max) {$\max\{x,0\}$};
    \draw[->] (xsplit) |- (max.west);

    % adder of max and xsq
    \node[circop,right=2cm of mul,fill=green!20] (add) {$+$};
    \draw[->] (mul) -- node[midway,above] {\texttt{xsq}} (add);
    \draw[->] (max.east) -| node[midway,above,xshift=0.3cm] {\texttt{r}} (add.south);
    \draw[->] (add) -- ++(1.5,0) node[right] {\texttt{y}};
\end{tikzpicture}
  \caption{Multi-path circuit.}
  \label{fig:multi_path}
\end{figure}

\subsection{Block diagram for a multiplexer}
\begin{systemverilogcode}
// mux delay = 1 ns
if (x > 0) begin  // delay = 0.5
  z = func1(x);  // delay = 4
end else begin 
  z = func2(x);  // delay = 2
end
y = func3(z);    // delay = 1.5
\end{systemverilogcode}

\begin{figure}
\centering

  \begin{tikzpicture}[node distance=1cm]

    % x, f1 and f2
    \node (x) {\texttt{x}};
    \coordinate[right=1cm of x] (xsplit);
    \node[draw, right=1cm of xsplit,minimum height=1cm,fill=green!10,yshift=1cm] 
      (f1) {$f_1(x)$};
    \node[draw, right=1cm of xsplit,minimum height=1cm,fill=green!10,yshift=-1cm] 
      (f2) {$f_2(x)$};  
    
    \draw[->] (x) -- (xsplit);
    \draw[->] (xsplit) |- (f1);   
    \draw[->] (xsplit) |- (f2);

    % mux block
    \node[draw, right=4cm of xsplit,minimum height=3cm,fill=green!10] (mux) {mux};
    \draw[->] (f1.east) -- (mux.west |- f1.east);
    \draw[->] (f2.east) -- (mux.west |- f2.east);

    % Selection block
    \node[draw, below=1cm of mux, minimum height=1cm, fill=green!10] (sel) {$x > 0$};
    \draw[->] (xsplit) |- (sel.west);
    \draw[->] (sel) -- (mux.south);

    % f3 block
    \node[draw, right=2cm of mux, minimum height=1cm, fill=green!10] (f3) {$f_3(z)$};
    \draw[->] (mux) -- node[midway,above] {\texttt{z}} (f3);
    \draw[->] (f3) -- ++(1.5,0) node[right] {\texttt{y}};

\end{tikzpicture}
\end{figure}



\begin{figure}
\centering

  \begin{tikzpicture}[node distance=1cm]

    % x, f1 and f2
    \node (x) {\texttt{x}};
    \coordinate[right=1cm of x] (xsplit);
    \node[draw, right=1cm of xsplit,minimum height=1cm,fill=green!10,yshift=1cm] 
      (f1) {$f_1(x)$};
    \node[draw, right=1cm of xsplit,minimum height=1cm,fill=green!10,yshift=-1cm] 
      (f2) {$f_2(x)$};  
    
    \draw[-,ultra thick, red] (x) -- (xsplit);
    \draw[->,ultra thick, red] (xsplit) |- (f1);   
    \draw[->] (xsplit) |- (f2);

    % mux block
    \node[draw, right=4cm of xsplit,minimum height=3cm,fill=green!10] (mux) {mux};
    \draw[->,ultra thick, red] (f1.east) -- (mux.west |- f1.east);
    \draw[->] (f2.east) -- (mux.west |- f2.east);

    % Selection block
    \node[draw, below=1cm of mux, minimum height=1cm, fill=green!10] (sel) {$x > 0$};
    \draw[->] (xsplit) |- (sel.west);
    \draw[->] (sel) -- (mux.south);

    % f3 block
    \node[draw, right=2cm of mux, minimum height=1cm, fill=green!10] (f3) {$f_3(z)$};
    \draw[->,ultra thick, red] (mux) -- node[midway,above] {\texttt{z}} (f3);
    \draw[->,ultra thick, red] (f3) -- ++(1.5,0) node[right] {\texttt{y}};

\end{tikzpicture}
\caption{Block diagram for a multiplexer highlight the critical path}
\end{figure}

\subsection{Register}

Consider the following SystemVerilog code for a register:
\begin{systemverilogcode}
  always_comb begin
    x_next = f(x);
  end
  always_ff @(posedge clk) begin
    x <= x_next;
  end
\end{systemverilogcode}


\begin{figure}
\begin{tikzpicture}[>=stealth, font=\scriptsize, node distance=1.8cm]

  % Register 
  \node[shape=register, draw, minimum width=1cm, minimum height=2cm,
        fill=blue!10] (reg) {};

  % Input, clock and output

  \node [left=1cm of reg.input] (D) {\texttt{D}};
  \node [left=1cm of reg.clk] (clk) {\texttt{clk}};
  \draw[->] (D) -- (reg.input);
  \draw[->] (clk) -- (reg.clk);
  \node [right=1cm of reg.output] (Q) {\texttt{Q}};
  \draw[->] (reg.output) -- (Q);

\end{tikzpicture}
\caption{Basic register }
\label{fig:reg_basic}
\end{figure}



\begin{figure}
\begin{tikzpicture}[>=stealth, font=\scriptsize, node distance=1.8cm]

  % Register 
  \node[shape=register, draw, minimum width=1cm, minimum height=2cm,
        fill=blue!10] (reg) {};

  % Input, clock and output
  \coordinate [left=1cm of reg.input] (xnext_in) {};
  \node [left=1cm of reg.clk] (clk) {\texttt{clk}};
  \draw[->] (xnext_in) -- (reg.input);
  \draw[->] (clk) -- (reg.clk);

  % Function block
  \node [draw, right=1cm of reg, minimum height=1cm, fill=green!10] (f) {$f(x)$};
  \draw[->] (reg.output) -- node[midway,above] {\texttt{x}} (f);
  \draw[-] (f.east) -- ++(1,0) -- ++(0,1.5) -| 
    node[midway,above,xshift=1cm] {\texttt{x\_next=f(x)}} (xnext_in);

\end{tikzpicture}
  \caption{Basic register block diagram of $x \leftarrow f(x)$}
  \label{fig:sync_circuit}
  
\end{figure}

\subsection{Two-Variable System}

Consider two-variable system:
\begin{systemverilogcode}
  always_comb begin
    z0 = func1(x0); 
    z1 = func2(x1);   
    x_next0 = func3(z0, z1);    
    x_next1 = func4(z0, x1);
  end
  always_ff @(posedge clk) begin
    x0 <= x_next0;
    x1 <= x_next1;
  end
\end{systemverilogcode}

Register-to-register paths are:
\begin{systemverilogcode}
x0 -> func1 -> func3 -> x0
x0 -> func3 -> x1
x1 -> func2 -> func3 -> x0
x1 -> func4 -> x1
\end{systemverilogcode}




\begin{figure}
\begin{tikzpicture}[>=stealth, font=\scriptsize, node distance=1.8cm]


  % Initial variables
  \node (x0) {$x_0$};
  \node [below=2cm of x0] (x1) {$x_1$};

  % Function blocks
  \node [draw, right=1cm of x0, minimum height=1cm, fill=green!10] (f1) {$f_1(x_0)$};
  \node [draw, below=0.2cm of f1, minimum height=1cm, fill=green!10] (f2) {$f_2(x_1)$};
  \node [draw, right=1cm of f1, minimum height=1cm, fill=green!10] (f3) {$f_3(z_0,z_1)$};
  \node [draw, right=2cm of x1, minimum height=1cm, fill=green!10] (f4) {$f_4(z_0,x_1)$};
  \draw[->] (x0) -- (f1);
  \draw[->] (x1) -- ++(0.5,0) |- (f2.west);
  \draw[->] (f1) -- node[midway,above] {$z_0$} (f3);
  \draw[->] (f2) -| node[midway,above,xshift=-0.5cm] {$z_1$} (f3.south);
  \draw[->] (f1.east) -- ++(0.5,0) -- ++(0,-0.5) -| (f4.north);;

  % Register 
  \node[shape=register, draw, minimum width=1cm, minimum height=2cm,
        right=6cm of x0,  fill=blue!10] (x0reg) {};
  \node[shape=register, draw, minimum width=1cm, minimum height=2cm,
        right=6cm of x1, fill=blue!10] (x1reg) {};
  \draw[->] (x1) -- (f4);
  \draw[->] (f3) -- node[midway,above] {\texttt{x\_next0}} (x0reg.input);
  \draw[->] (f4) -- node[midway,above] {\texttt{x\_next1}} (x1reg.input);
  \draw[->] (x0reg.output) -- ++(0.5,0) node [right] {$x_0'$};
  \draw[->] (x1reg.output) -- ++(0.5,0) node [right] {$x_1'$};
\end{tikzpicture}
  \caption{Simple two variable system}
  \label{fig:twovar}
  
\end{figure}



\begin{figure}
\begin{tikzpicture}[>=stealth, font=\scriptsize, node distance=1.8cm]


  % Initial variables
  \node (x0) {$x_0$};
  \node [below=2cm of x0] (x1) {$x_1$};

  % Function blocks
  \node [draw, right=1cm of x0, minimum height=1cm, fill=green!10] (f1) {$f_1(x_0)$};
  \node [draw, below=0.2cm of f1, minimum height=1cm, fill=green!10] (f2) {$f_2(x_1)$};
  \node [draw, right=1cm of f1, minimum height=1cm, fill=green!10] (f3) {$f_3(z_0,z_1)$};
  \node [draw, right=2cm of x1, minimum height=1cm, fill=green!10] (f4) {$f_4(z_0,x_1)$};
  \draw[->] (x0) -- (f1);
  \draw[->,ultra thick, red] (x1) -- ++(0.5,0) |- (f2.west);
  \draw[->] (f1) -- node[midway,above] {$z_0$} (f3);
  \draw[->,ultra thick, red] (f2) -| node[midway,above,xshift=-0.5cm] {$z_1$} (f3.south);
  \draw[->] (f1.east) -- ++(0.5,0) -- ++(0,-0.5) -| (f4.north);;

  % Register 
  \node[shape=register, draw, minimum width=1cm, minimum height=2cm,
        right=6cm of x0,  fill=blue!10] (x0reg) {};
  \node[shape=register, draw, minimum width=1cm, minimum height=2cm,
        right=6cm of x1, fill=blue!10] (x1reg) {};
  \draw[->] (x1) -- (f4);
  \draw[->,ultra thick, red] (f3) -- node[midway,above] {\texttt{x\_next0}} (x0reg.input);
  \draw[->] (f4) -- node[midway,above] {\texttt{x\_next1}} (x1reg.input);
  \draw[->,ultra thick, red] (x0reg.output) -- ++(0.5,0) node [right] {$x_0'$};
  \draw[->] (x1reg.output) -- ++(0.5,0) node [right] {$x_1'$};
\end{tikzpicture}
  \caption{Critical path in two variable system}
  \label{fig:twovar_critical}
  
\end{figure}




\begin{figure}
\begin{tikzpicture}[>=stealth, font=\scriptsize, node distance=1.8cm]

  % Register 
  \node[shape=register, draw, minimum width=1cm, minimum height=2cm,
        fill=blue!10] (reg) {};

  % Input, clock and output
  \coordinate [left=1cm of reg.input] (xnext_in) {};
  \node [left=1cm of reg.clk] (clk) {\texttt{clk}};
  \draw[->] (xnext_in) -- (reg.input);
  \draw[->] (clk) -- (reg.clk);

  % Function block
  \node [draw, right=1cm of reg, minimum height=1cm, fill=green!10] (f) {$f(x)$};
  \draw[->] (reg.output) -- node[midway,above] {\texttt{x}} (f);
  \draw[-] (f.east) -- ++(1,0) -- ++(0,1.5) -| 
    node[midway,above,xshift=1cm] {\texttt{x\_next=f(x)}} (xnext_in);

\end{tikzpicture}
  \caption{Basic register block diagram of $x \leftarrow f(x)$}
  \label{fig:sync_circuit}
  
\end{figure}

\subsection{Minimum Hold Time}

Consider the following SystemVerilog code:
\begin{systemverilogcode}
always_ff @(posedge clk) begin
  x0 <= x;
  x1 <= x0;
end
\end{systemverilogcode}

This is a simple delay line:
\begin{figure}
\begin{tikzpicture}[>=stealth, font=\scriptsize, node distance=1.8cm]

  % Initial variables
  \node (x) {$x$};

  % Register blocks
  \node[shape=register, draw, minimum width=1cm, minimum height=2cm,
        right=1cm of x, fill=blue!10] (reg0) {};
  \node[shape=register, draw, minimum width=1cm, minimum height=2cm,
        right=1cm of reg0, fill=blue!10] (reg1) {};

  % Connections
  \draw[->] (x) -- (reg0.input);
  \draw[->] (reg0.output) -- node[midway,above] {\texttt{x0}} (reg1.input);
  \draw[->] (reg1.output) -- ++(0.5,0) node [right] {\texttt{x1}};

\end{tikzpicture}
  \caption{Simple delay line}
  \label{fig:delay_line}
  
\end{figure}
\end{document}
